\subsection{A modular nondistributive dominion lattice: Problem 16.20}
\label{cand:16.20}
\problemstatement{16.20}

For \(H\leq A\) and a class \(\mathcal N\) of target groups,
the dominion consists of those \(a\in A\) on which every pair
of homomorphisms from \(A\) into a member of \(\mathcal N\)
agree whenever they agree on \(H\). The literature convention
allows \(A\notin\mathcal N\); see
\cite{shakhova-dominions}. Here is the smaller of the archive's
two constructions under that convention.

Let
\[
 U_1=\SL_2(5),\quad U_2=\SL_2(7),\quad U_3=\SL_2(13),
 \qquad A=(U_1\times U_2\times U_3)/\langle z_1z_2z_3\rangle,
\]
where \(z_i=-I\), and take \(H=1\).
These standard perfect central covers have simple quotients
\(S_i=\operatorname{PSL}_2(p_i)\) of orders \(60,168,1092\).
No one of these three orders divides another. The center
\(W=Z(A)\cong C_2^2\) has as its three nonzero vectors the
images \(v_i\) of the \(z_i\). The group has order
\(44{,}029{,}440\).

Every normal subgroup \(R\) of \(A\) is \(A_I K\), where
\(A_I\) is the product of the factors indexed by some
\(I\subseteq\{1,2,3\}\), and
\[
 K=R\cap W,\qquad \langle v_i:i\in I\rangle\leq K\leq W.
\]
To see this, its image in the product of the three simple
quotients is a product of factors. For any occurring factor,
\([R,U_i]\) projects onto \(S_i\). A normal subgroup of
\(U_i\) projecting onto \(S_i\) is all of \(U_i\):
if \(U_i=JZ(U_i)\), perfectness gives \(U_i\leq J\).
Removing those factors leaves precisely a central subgroup.

The key fact is that every homomorphism \(f:A\to A/R\)
kills \(R\). Project its restriction on \(U_i\) to each
surviving \(S_j\). For \(j\ne i\) that map is trivial.
Indeed, every proper normal subgroup of \(U_i\) is central,
so every nontrivial image has order divisible by \(|S_i|\),
which cannot divide \(|S_j|\). Thus a removed factor has
central image and is killed by perfectness. A surviving
factor maps into its own image \(D_i\) times the center;
perfectness puts its image in
\([D_iZ(A/R),D_iZ(A/R)]=D_i\).
When \(v_i\in K\), this \(D_i\) is \(S_i\), and every
homomorphism \(U_i\to S_i\) kills \(z_i\). Every nonzero
vector of \(K\) is some \(v_i\), so \(f\) kills \(K\)
as well as the removed factors, proving the fact.

Let \(q(\mathcal C)\) denote the quasivariety generated by
\(\mathcal C\), and put
\(\mathcal M=q(\{A/R:R\trianglelefteq A\})\).
For a finite domain \(A\), the condition defining membership
of an element in a dominion is a finite quasiidentity in the
target: use two variables for each element of \(A\), impose
both multiplication tables and agreement on \(H\), and
conclude agreement at that element. It follows that the
dominion for \(q(\mathcal C)\) is exactly the intersection
of the dominions for the generating targets.

For \(H=1\), a dominion is simply an intersection of kernels
of maps into the target class, by comparison with the trivial
map. It is therefore normal. For each \(R\), all maps into
\(A/R\) kill \(R\), while the canonical quotient has kernel
exactly \(R\). The preceding transfer gives
\[
 \operatorname{dom}^{A}_{q(A/R)}(1)=R.
\]
Every normal subgroup occurs, and the entire dominion family
is the normal subgroup lattice of \(A\). That lattice is
modular by the Dedekind identity, and its interval
\([1,W]\) is the five-element diamond \(M_3\), so it is
not distributive.

The larger archived example also handles the extra convention
that every admissible target quasivariety contain \(A\).
It uses seven covers \(\SL_2(p)\),
\(p=7,11,13,17,19,23,29\), with their central involutions
mapped onto all nonzero vectors of \(\F_2^3\), and a fixed
central line \(H\). Under that convention the dominions are
exactly the five subspaces between \(H\) and \(\F_2^3\).
The smaller construction above suffices for the documented
convention of the printed question. Source:
\artifact{research/16.20-proof.md}.

